\documentclass[12pt]{article}
\usepackage[utf8]{inputenc}
\usepackage[T1]{fontenc}
\usepackage[margin=1in]{geometry}
\usepackage{setspace}
\usepackage{xcolor}
\usepackage[colorlinks=true, urlcolor=blue!60!black, linkcolor=blue!60!black]{hyperref}
\setlength{\parindent}{0em}
\setlength{\parskip}{0.45em}
\renewcommand{\baselinestretch}{1.0}
\begin{document}
\begin{center}
{\Large \textit{South Dakota v. Dole}, 483 U.S. 203 (1987)\par}
\vspace{0.3em}
{\normalsize Case Brief --- One-Page Edition\par}
\end{center}

\textbf{Memory Jogger.} The drinking-age highway-funds case: the five-part test for spending conditions (\href{https://supreme.justia.com/cases/federal/us/483/203/}{483 U.S. 203}).

\textbf{Facts.} Congress withheld 5\% of specified highway funds from States allowing people under 21 to buy alcohol; South Dakota let 19-year-olds buy beer containing up to 3.2\% alcohol. Congress acted after a presidential commission found that inconsistent state drinking ages encouraged young adults to cross state lines to drink and drive.

\textbf{Procedural History.} South Dakota lost in the District Court and the Eighth Circuit.

\textbf{Votes.} Affirmed, 7--2. Rehnquist wrote for the Court; Brennan and O'Connor dissented.

\textbf{Holding.} Congress may attach conditions to federal funds that serve the general welfare, are unambiguous, relate to the federal interest, induce nothing unconstitutional, and are not coercive, even to reach goals it cannot regulate directly.

\textbf{Reasoning.} The spending power reaches beyond Article I's enumerated fields. The condition related to safe interstate travel, raising the drinking age was lawful, and 5\% of highway funds was ``relatively mild encouragement.'' O'Connor called the condition a regulation of who may drink, not of how highway money is spent.

\textbf{Examples.} \emph{Same-Side:} a 0.08 blood-alcohol condition. \emph{Opposite-Side:} all K--12 funds for a campus gun ban (\textit{NFIB}). \emph{Fence-Sitter:} 25\% of education funds for a phone ban.

\textbf{Critique.} A ``Trojan Horse'' for federal power, with no bright line for coercion.

\textbf{Key Quotation.} ``These cases establish that the `independent constitutional bar' limitation on the spending power is not, as petitioner suggests, a prohibition on the indirect achievement of objectives which Congress is not empowered to achieve directly. Instead, we think that the language in our earlier opinions stands for the unexceptionable proposition that the power may not be used to induce the States to engage in activities that would themselves be unconstitutional.'' --- Chief Justice Rehnquist, Opinion of the Court, 483 U.S. at 210.

\end{document}
